%%%%%%%%%%%%%%%%%%%%%%%%%%%%%%%%%%%%%%%%%%%%%%%%%%%%%%%%%%%%%%%%%%%%%%
% Academic CV -- Subrat Prasad Panda (postdoc / faculty applications)
% Same visual system as cv_industry.tex (Cochineal + Cabin, green swish
% headings, crimson keys and bullets).
% Publications are generated from references.bib.
% Build: pdflatex -> biber -> pdflatex -> pdflatex
%%%%%%%%%%%%%%%%%%%%%%%%%%%%%%%%%%%%%%%%%%%%%%%%%%%%%%%%%%%%%%%%%%%%%%
\documentclass[11pt]{article}

\usepackage[a4paper,hmargin=1.8cm,top=1.0cm,bottom=1.1cm]{geometry}
\usepackage[T1]{fontenc}
\usepackage[utf8]{inputenc}
\usepackage[english]{babel}
\usepackage[p,lf]{cochineal}
\usepackage[medium,bold]{cabin}
\usepackage[varqu,varl,scale=0.88]{zi4}
\usepackage{amssymb}
\usepackage{enumitem}
\usepackage[dvipsnames]{xcolor}
\usepackage{tikz}
\usepackage{calc}
\usepackage{ragged2e}
\setlength{\JustifyingParindent}{0pt}
\usepackage{needspace}
\usepackage{fontawesome5}
\usepackage{simpleicons}
\usepackage{accsupp}
\usepackage[hyphens]{url}
\usepackage[style=ieee,sorting=ydnt,maxnames=99,defernumbers=true]{biblatex}
\usepackage{csquotes}
\usepackage[colorlinks=true,allcolors=black,urlcolor=Marker!80!black]{hyperref}

\input{glyphtounicode}
\pdfgentounicode=1

%------------------------------- PALETTE ------------------------------
\definecolor{Swish}{HTML}{88AC0B}
\definecolor{Marker}{HTML}{B6073F}
\definecolor{Ink}{HTML}{222222}
\definecolor{Muted}{HTML}{5F6368}
\definecolor{Sub}{HTML}{1F5A96}  % blue for subsection headings
\color{Ink}

\pagestyle{empty}
\setlength{\parindent}{0pt}
\raggedright
\urlstyle{same}

%------------------------------- HELPERS ------------------------------
% Icon that extracts as text (#2) in the PDF text layer
\newcommand{\icon}[2]{\BeginAccSupp{method=escape,ActualText={#2}}%
  \makebox[1.3em]{\color{Marker!85!black}#1}\EndAccSupp{}\hspace{1pt}}

% Section heading with gradient swish
\newcommand{\cvsection}[1]{%
  \par\needspace{4\baselineskip}\vspace{13pt}%
  \begin{tikzpicture}[baseline]
    \shade[left color=Swish!60!white,right color=white] (0,0) rectangle (\linewidth,2.5pt);
    \node[font=\Large\bfseries\sffamily,inner sep=0pt,anchor=south west,
          text depth=.5ex,text height=1.5ex] at (1pt,2pt) {#1};
  \end{tikzpicture}\par\vspace{4pt}}

% Left key column (dates / labels)
\newlength{\keywidth}\setlength{\keywidth}{2.3cm}
\newcommand{\keyfont}{\sffamily\footnotesize\bfseries\color{Marker!80!black}}
\newcommand{\keylabel}[1]{\raisebox{0pt}[\height][0pt]{\parbox[t]{\labelwidth}{\leavevmode\raggedright\keyfont #1}}}

% \begin{cventry}{key}{Title}{Right-aligned info}{Subtitle line (optional)}
\newenvironment{cventry}[4]{%
  \begin{list}{}{\setlength{\leftmargin}{\keywidth}\setlength{\labelwidth}{\keywidth-0.4em}%
    \setlength{\labelsep}{0.4em}\setlength{\topsep}{0pt}\setlength{\itemsep}{0pt}%
    \setlength{\parsep}{0pt}\setlength{\partopsep}{0pt}}
  \item[{\keylabel{#1}}]%
    {\sffamily\bfseries #2}\hfill{\sffamily\small\color{Muted}#3}%
    \ifx\relax#4\relax\else\\{\small\itshape\color{Muted}#4}\fi\par
}{\end{list}\vspace{5pt}}

% Single key/value row (gap after each row adjustable per section)
\newlength{\rowgap}\setlength{\rowgap}{2.5pt}
\newcommand{\cvrow}[2]{%
  \begin{list}{}{\setlength{\leftmargin}{\keywidth}\setlength{\labelwidth}{\keywidth-0.4em}%
    \setlength{\labelsep}{0.4em}\setlength{\topsep}{0pt}\setlength{\itemsep}{0pt}\setlength{\parsep}{0pt}\setlength{\partopsep}{0pt}}
  \item[{\keylabel{#1}}] #2
  \end{list}\vspace{\rowgap}}

% Blue subsection heading inside a section
\newcommand{\subgroup}[1]{{\sffamily\bfseries\color{Sub}#1}\par\vspace{3pt}}

\setlist[itemize]{leftmargin=1.05em,itemsep=1pt,topsep=2pt,parsep=0pt,
  label={\color{Marker!80!black}\scriptsize$\blacktriangleright$}}

%------------------------------- PUBLICATIONS -------------------------
\addbibresource{references.bib}

% Bold own name
\DeclareNameFormat{given-family}{%
  \ifboolexpr{test {\ifdefstring{\namepartfamily}{Panda}}}
    {\textbf{\usebibmacro{name:given-family}{\namepartfamily}{\namepartgiveni}{\namepartprefix}{\namepartsuffix}}}
    {\usebibmacro{name:given-family}{\namepartfamily}{\namepartgiveni}{\namepartprefix}{\namepartsuffix}}%
  \usebibmacro{name:andothers}}

% Numbered crimson circles
\newcommand*\circled[1]{\tikz[baseline=(char.base)]{
  \node[shape=circle,text=white,fill=Marker!80!black,font=\sffamily\scriptsize\bfseries,
        inner sep=1pt,text height=1.35ex,minimum width=1.5em,text centered] (char) {#1};}}
\newcounter{bibitem}
\AtBeginBibliography{\setcounter{bibitem}{1}}
\DeclareFieldFormat{labelnumberwidth}{\makebox[2.4em][l]{\circled{\thebibitem}}}
\AtEveryBibitem{\stepcounter{bibitem}}
\setlength{\biblabelsep}{0pt}

% Compact entries; titles kept exactly as in the .bib, set in semibold sans
\renewcommand*{\bibfont}{\small}
\setlength{\bibitemsep}{0.9ex}
\setlength{\bibhang}{2.4em}
\DeclareFieldFormat{titlecase}{#1}
\DeclareFieldFormat{sentencecase}{#1}
\DeclareFieldFormat*{title}{{\sffamily\bfseries #1}\isdot}
\AtEveryBibitem{\clearfield{url}\clearfield{doi}\clearfield{month}\clearlist{organization}}

\defbibheading{pubgroup}{\par\vspace{5pt}{\sffamily\bfseries\color{Sub}#1}\par\vspace{3pt}}

% references.bib uses website-friendly entry types. Map them onto
% biblatex types, remembering the original kind in `entrysubtype':
%   @journal -> article        @conference -> inproceedings
%   @workshop -> misc/workshop @talk -> misc/talk
%   @underreview -> misc/underreview (no venue: double-blind)
% New entries only need the right @type; no edits here.
\DeclareSourcemap{
  \maps[datatype=bibtex]{
    \map{\step[typesource=journal,     typetarget=article,       final]}
    \map{\step[typesource=conference,  typetarget=inproceedings, final]}
    \map{\step[typesource=workshop,    typetarget=misc,          final]
         \step[fieldset=entrysubtype,  fieldvalue=workshop]}
    \map{\step[typesource=talk,        typetarget=misc,          final]
         \step[fieldset=entrysubtype,  fieldvalue=talk]}
    \map{\step[typesource=underreview, typetarget=misc,          final]
         \step[fieldset=entrysubtype,  fieldvalue=underreview]
         \step[fieldset=howpublished,  fieldvalue={Under review}]}
  }
}

% Topic chips from the `keywords' field, shown after each entry
\newcommand{\pubchip}[1]{\tikz[baseline=(c.base)]\node[fill=Sub!8,draw=Sub!35,line width=0.4pt,
  rounded corners=1.5pt,inner xsep=2.5pt,inner ysep=0.8pt,text=Sub,
  font=\sffamily\scriptsize] (c) {#1};\hspace{2pt}}
\renewbibmacro*{finentry}{\finentry
  \iffieldundef{keywords}{}{\hspace{0.4em}\forcsvfield{\pubchip}{keywords}}}

\begin{document}

%================================ HEADER ===============================
% \contact{icon}{text for parsers}{content}
\newcommand{\contact}[3]{\mbox{\icon{#1}{#2}#3}}
{\fontsize{25}{28}\selectfont\sffamily\bfseries Subrat Prasad Panda}\\[3pt]
{\sffamily\large\color{Marker!85!black} Ph.D. Candidate, Computer Science, NTU Singapore}\\[1pt]
{\sffamily\small\color{Muted} Trustworthy RL \textperiodcentered{} Neurosymbolic AI}\\[6pt]
{\sffamily\small
\contact{\faEnvelope[regular]}{Email: }{\href{mailto:subratpr001@e.ntu.edu.sg}{subratpr001@e.ntu.edu.sg}}\hfill
\contact{\faMobile*}{Phone: }{+65 ****3021}\hfill
\contact{\faGlobe}{Website: }{\href{https://subratpp.github.io/}{subratpp.github.io}}\hfill
\contact{\faMapMarker*}{Location: }{Singapore}\\[2pt]
\contact{{\scriptsize\simpleicon{googlescholar}}}{Scholar: }{\href{https://scholar.google.com/citations?user=-kOZ8c0AAAAJ&hl=en}{Google Scholar}}\hfill
\contact{\faGithub}{GitHub: }{\href{https://github.com/subratpp}{github.com/subratpp}}\hfill
\contact{\faLinkedin}{LinkedIn: }{\href{https://www.linkedin.com/in/subratpp/}{linkedin.com/in/subratpp}}\par}
\vspace{4pt}


%================================ RESEARCH =============================
\cvsection{Research Statement}
{\justifying
I am broadly interested in \textbf{Trustworthy Reinforcement Learning}, particularly in designing interpretable, robust, and verifiable policies for safety-critical domains. My current work focuses on \textbf{Neurosymbolic} approaches such as programmatic policies, combining explicit knowledge with learning to bridge System~1 (neural) and System~2 (symbolic) reasoning.\par}

%================================ SKILLS ===============================
\cvsection{Research Areas and Skills}
{\setlength{\rowgap}{5.5pt}
\cvrow{Reinforcement Learning}{Hierarchical and goal-conditioned RL, sample-efficient learning in sparse-reward and long-horizon tasks, robust and distributional RL, RL fine-tuning of foundation models}
\cvrow{Neurosymbolic AI}{Programmatic and interpretable policies (differentiable decision trees), symbolic planning with incremental knowledge, interpretable state abstractions and world models}
\cvrow{Safety \& Verification}{Formal verification of learned controllers, reach-avoid and safety certificates, abstraction-based reasoning for safety-critical systems}
\cvrow{Applications}{Cyber-physical systems, microgrid energy management, autonomous navigation and robotics, LLM-based planning for embodied agents}
\cvrow{Programming}{Python, C, PyTorch, TensorFlow, ROS 2, MATLAB, \LaTeX}
\cvrow{Coursework}{Mathematical Statistics, Deep Learning, Machine Learning, Computer Vision}
}

%================================ EDUCATION ============================
\cvsection{Education}

\begin{cventry}{2023 -- 2027\\[-1pt]{\mdseries\color{Muted}expected}}{Ph.D., Computer Science}{Singapore}{Nanyang Technological University (NTU)}
Thesis: \emph{Neurosymbolic Techniques for Safe and Efficient Reinforcement Learning}\\
Advisors: Arvind Easwaran (NTU), Blaise Genest (IPAL, CNRS France \& Singapore)
\end{cventry}

\begin{cventry}{2019 -- 2021}{M.Tech., Computer Science}{Kolkata, India}{Indian Statistical Institute (ISI)}
Thesis: \emph{Strategies for User Allocation and Service Placement in Multi-Access Edge Computing}\\
Advisor: Ansuman Banerjee \quad CPI: 82.1\%
\end{cventry}

\begin{cventry}{2012 -- 2016}{B.Tech., Electronics and Instrumentation}{Silchar, India}{National Institute of Technology (NIT) Silchar}
Thesis: \emph{Real-Time Position Control of DC Motor using PID and Fuzzy Logic}\\
Advisor: Munmun Khanra \quad CPI: 7.92/10
\end{cventry}

%================================ EXPERIENCE ===========================
\cvsection{Research Experience}

\begin{cventry}{2022 -- 2023}{Research Associate}{Singapore}{CNRS@CREATE}
\begin{itemize}
  \item Designed a granular control system for microgrids using Deep Reinforcement Learning.
\end{itemize}
\end{cventry}

\begin{cventry}{2021 -- 2022}{Machine Learning Engineer}{Pune, India}{InFocusp Innovation Pvt.\ Ltd., under attachment to Google X (USA)}
\begin{itemize}
  \item Developed solutions for the Abstraction and Reasoning Challenge (ARC) using Reinforcement Learning.
\end{itemize}
\end{cventry}

\begin{cventry}{2021 -- 2021}{Research Engineer}{Hyderabad, India}{TiHAN, IIT Hyderabad}
\begin{itemize}
  \item Developed braking and throttle control systems using Arduino for a test vehicle.
  \item Deployed a lane detection model on the NVIDIA AGX board.
\end{itemize}
\end{cventry}

%================================ PUBLICATIONS =========================
\cvsection{Publications}
\nocite{*}
\printbibliography[heading=pubgroup,title={Under Review},subtype=underreview]
\printbibliography[heading=pubgroup,title={Journal Articles},type=article]
\printbibliography[heading=pubgroup,title={Conference Proceedings},type=inproceedings]
\printbibliography[heading=pubgroup,title={Workshop Papers},subtype=workshop]
\printbibliography[heading=pubgroup,title={Talks and Tutorials},subtype=talk]

%================================ SOFTWARE =============================
\cvsection{Open-Source Software and Educational Content}
% Name links to the repository
\newcommand{\repo}[2]{\href{#2}{#1\,\BeginAccSupp{method=escape,ActualText={}}{\scriptsize\faGithub}\EndAccSupp{}}}

\subgroup{Research Code}
\cvrow{\repo{DTSemNet}{https://github.com/CPS-research-group/dtsemnet}}{Differentiable oblique decision trees and tree policies for RL (JMLR 2026, ECAI 2024).}
\cvrow{\repo{InK-BWTS}{https://github.com/CPS-research-group/ink_bwts}}{Neurosymbolic hierarchical RL with Belief World Tree Search (ECML-PKDD 2026).}
\cvrow{\repo{F1Tenth}{https://github.com/subratpp/f1tenth}}{ROS 2 race stack (SLAM, particle-filter localisation, Pure Pursuit, etc.) on Jetson AGX.}

\vspace{3pt}
\subgroup{Educational Content {\mdseries\small\color{Muted}(ongoing effort)}}
\cvrow{\repo{Deep RL Jutsu}{https://github.com/subratpp/deep-rl-jutsu}}{Notebook-style tutorials implementing deep RL algorithms from scratch (DQN, DDPG, SAC, PPO), with a curated guide to RL libraries and resources.}
\cvrow{\repo{GenAI}{https://github.com/subratpp/generative_ai}}{Small-scale implementations of generative models (VAE, GAN, diffusion, state-space models, CLIP, VLA) on toy problems to explain how they work.}
\cvrow{\repo{RLFT}{https://github.com/subratpp/rlft}}{Tutorial series on RL fine-tuning of foundation models: RLHF, DPO and GRPO, from a from-scratch GPT to a pretrained LLM.}

%================================ AWARDS ===============================
\cvsection{Awards and Scholarships}
\cvrow{2023 -- 2027}{\textbf{DesCartes Scholarship}: Ph.D. scholarship awarded by CNRS@CREATE under the DesCartes project (research collaboration between Singapore and France).}
\cvrow{2021}{\textbf{M.Tech Best Thesis Nomination}: ranked among the top 5 theses of the batch.}
\cvrow{2019 -- 2021}{\textbf{M.Tech Scholarship}: scholarship awarded to pursue a master's degree.}
\cvrow{2006 -- 2012}{\textbf{Jawahar Navodaya Vidyalaya (JNV) Scholarship}: fully funded residential education from Class 6 to 12.}

%================================ ACADEMIC SERVICES ====================
\cvsection{Academic Services}

\subgroup{Teaching and Mentoring}
\begin{cventry}{Final Year\\Projects}{Undergraduate Final Year Projects}{NTU}{Project supervision}
\emph{SafeDAG: Precedence-Constrained LLM Steering for Safer Planning in Embodied Agents} (ongoing); programmatic policies for interpretable RL using pre-trained models / LLMs; robustness vs.\ explainability trade-offs in CNNs.
\end{cventry}

\begin{cventry}{URECA}{Undergraduate Research Experience on Campus}{NTU}{Research mentoring}
Mechanistic interpretability of Decision Transformers; ARC reasoning with LLMs; physics-inspired model-based RL; residual reinforcement learning.
\end{cventry}

\begin{cventry}{F1Tenth}{F1Tenth Autonomous Racing Team}{NTU}{Mentor}
Built the ROS 2 race stack on Jetson AGX and advised students on control and navigation strategies.
\end{cventry}

\subgroup{Reviewing}
\cvrow{Reviewer}{CAV, AAAI, ECAI, IEEE TSMC, RTSS, NeurIPS, AAMAS, NeSy.}

\vspace{3pt}
\subgroup{Leadership}
\cvrow{Vice President}{\textbf{CCDS Student Club}, NTU: organized events such as the Deep Learning Bootcamp and AI Meetup.}

%================================ REFERENCES ===========================
\needspace{6\baselineskip}
\cvsection{Academic References}
% \referee{relationship}{Name}{Position, Affiliation}{email (may be empty)}
\newcommand{\referee}[4]{\cvrow{#1}{\textbf{#2}, #3\ifx\relax#4\relax\else\hfill{\small\href{mailto:#4}{#4}}\fi}}
\referee{Ph.D. Advisor}{Arvind Easwaran}{Associate Professor, CCDS, NTU Singapore}{arvinde@ntu.edu.sg}
\referee{Ph.D. Advisor}{Blaise Genest}{Research Scientist, CNRS@CREATE, Singapore}{blaise.genest@cnrsatcreate.sg}
\referee{M.Tech Advisor}{Ansuman Banerjee}{Professor, ISI Kolkata, India}{ansuman@isical.ac.in}

%================================ COLLABORATORS ========================
\cvsection{Academic Collaborators}
\referee{Verification}{\DJ{}or\dj{}e \v{Z}ikeli\'c}{Assistant Professor, CCDS, NTU Singapore}{djordje.zikelic@ntu.edu.sg}
\referee{Decision Trees}{P. N. Suganthan}{Research Professor, KINDI, Qatar University}{p.n.suganthan@qu.edu.qa}
\referee{Edge AI}{Arani Bhattacharya}{Assistant Professor, IIIT Delhi, India}{arani@iiitd.ac.in}


\end{document}
